В данной главе приводится анализ результатов разработки криптографической библиотеки. Оценивается эффективность принятых архитектурных решений, таких как использование избыточного представления Radix-51 и применение векторизации. Также обсуждаются результаты профилирования, сравнение производительности с библиотекой Bouncy Castle и определяются возможные направления развития проекта.

\subsection{Оценка эффективности разработанной криптографической библиотеки}

Оценка эффективности позволяет проверить, насколько итоговый программный продукт соответствует заявленным целям исследования и техническим требованиям.

\subsubsection{Сопоставление полученных результатов с поставленными задачами}

В начале исследования были поставлены следующие задачи:
\begin{enumerate}
    \item Изучить математический аппарат кривых Монтгомери и Эдвардса, а также специфику вычислений в конечном поле $GF(2^{255}-19)$.
    \item Спроектировать вычислительное ядро библиотеки на базе представления Radix-51.
    \item Разработать программный каркас, исключающий выделение памяти в куче (Zero Allocation).
    \item Реализовать алгоритмы пакетной (batched) обработки с использованием векторных инструкций AVX2.
    \item Сравнить производительность разработанного решения с библиотекой Bouncy Castle.
\end{enumerate}

В соответствии с поставленными задачами были получены следующие практические результаты:
\begin{enumerate}
    \item Изучена теоретическая база протоколов X25519 и Ed25519. В коде реализованы полные законы сложения и работа в проективных координатах, что исключило необходимость использования ветвлений и обеспечило константное время выполнения операций (constant-time).
    
    \item Реализовано вычислительное ядро на базе Radix-51. Элементы поля представлены в виде пяти 64-битных регистров, что минимизировало количество операций нормализации переносов по сравнению с классическим представлением Radix-25.5.
    
    \item Разработана архитектура управления памятью. Благодаря использованию типов ref struct и Span<byte>, криптографические операции выполняются без аллокации объектов в управляемой куче.
    
    \item Внедрена пакетная обработка. Произведен переход от структур данных AoS (Array of Structures) к SoA (Structure of Arrays). С помощью интринсиков System.Runtime.Intrinsics реализована возможность одновременного вычисления четырех независимых скалярных умножений.
    
    \item Выполнено тестирование и профилирование. Библиотека успешно проходит тесты по векторам из RFC 7748 и RFC 8032. Профилирование в BenchmarkDotNet показало, что при пакетной обработке библиотека опережает Bouncy Castle по пропускной способности в среднем на 30–40\%.
\end{enumerate}

Таким образом, поставленные в начале работы задачи выполнены в полном объеме. Главная цель — создание высокопроизводительной криптографической библиотеки X25519/Ed25519 для платформы .NET — достигнута.

\subsubsection{Анализ преимуществ и недостатков реализованного решения}

На основе этапов разработки и тестирования выделены следующие объективные сильные и слабые стороны библиотеки.

Преимущества:
\begin{enumerate}
    \item Улучшенная производительность. Векторизация операций (AVX2) и применение Radix-51 позволили добиться более высокой скорости вычислений по сравнению со скалярными реализациями на C\#, такими как Bouncy Castle.
    \item Полностью управляемый код (Managed Code). Библиотека написана на C\# без использования P/Invoke и вызовов нативных C/C++ библиотек. Это упрощает сборку и развертывание проекта на различных операционных системах (Windows, Linux, macOS).
    \item Снижение нагрузки на сборщик мусора. Отсутствие аллокаций (Zero Allocation) в критических участках кода решает проблему скачков задержек (latency spikes), характерную для высоконагруженных .NET-приложений.
    \item Безопасность выполнения. Использование побитовых масок вместо операторов условного перехода (CSWAP) защищает алгоритм от базовых атак по сторонним каналам (timing attacks).
\end{enumerate}

Недостатки и ограничения:
\begin{enumerate}
    \item Аппаратная зависимость максимальной скорости. Пиковая пропускная способность при пакетной обработке достигается только на процессорах с поддержкой набора инструкций AVX2 (архитектура x86/x64). 
    \item Отсутствие векторизации для ARM. В текущей версии не реализовано использование интринсиков NEON. На серверах с ARM-процессорами библиотека переключается в скалярный режим работы (fallback), что снижает её эффективность.
    \item Узкая специализация. Библиотека реализует только стандарты на базе кривой Curve25519. Поддержка других распространенных стандартов (например, алгоритмов на кривых NIST) не предусмотрена.
\end{enumerate}

Перечисленные ограничения характерны для специализированных криптографических модулей первой версии и оставляют пространство для дальнейшей оптимизации программного кода.

\subsection{Обсуждение результатов бенчмаркинга и профилирования}

Результаты тестирования в BenchmarkDotNet подтвердили гипотезу о том, что архитектура SoA и применение 256-битных векторных инструкций способны нивелировать накладные расходы управляемой среды выполнения C\#. 

При пакетной обработке (обновление четырех ключей одновременно) среднее время, затрачиваемое процессором на одну операцию, оказалось заметно ниже показателей Bouncy Castle. 
Прирост пропускной способности составил около 30–40\% в зависимости от конкретного криптографического примитива (X25519 или Ed25519).

Важный результат показало профилирование памяти. В тестах Bouncy Castle наблюдается линейный рост объема выделяемой памяти (до сотен мегабайт при длительных нагрузках), что приводит к частым вызовам сборщика мусора (Gen0 Collections).
В разработанной библиотеке метрика Allocated остается стабильной и близкой к нулю независимо от количества итераций, а сборки мусора в процессе выполнения криптографических расчетов полностью отсутствуют. Это делает решение оптимальным для серверов, требующих стабильного времени отклика.

\subsection{Определение направлений для дальнейшего развития продукта}

Исходя из архитектуры проекта и выявленных ограничений, можно определить следующие направления для доработки и развития библиотеки:

\begin{enumerate}
    \item Поддержка AVX-512. Использование новых 512-битных инструкций процессора (доступных в последних версиях .NET) позволит увеличить размер пакета с 4 до 8 операций за один такт, что обеспечит дополнительное масштабирование производительности на современных серверных процессорах.
    \item Внедрение интринсиков для ARM (AdvSimd). Добавление реализации пакетной обработки на базе 128-битных регистров для архитектуры ARM64 позволит задействовать векторное ускорение на облачных платформах (например, AWS Graviton или Apple Silicon).
    \item Интеграция со стандартными интерфейсами CLR. Написание классов-адаптеров для абстракций System.Security.Cryptography.AsymmetricAlgorithm упростит внедрение разработанной библиотеки в уже существующие корпоративные проекты.
\end{enumerate}